%% University of Edinburgh School of Mathematics thesis/report starter —
%% edmaths style.
%%
%% edmaths is NOT included in this download: edmaths.sty ships with TeX Live
%% (and with latextemplates.com's compiler), so bundling a copy would only pin
%% you to an older version than your own distribution already has.
%%
%% Documentation: texdoc edmaths  |  https://ctan.org/pkg/edmaths
%% Licence: edmaths.sty is LPPL-1.3-or-later, (c) 2025 Josh Fogg
%% (originally written by Thomas Koeppe, 2007).
%%
%% edmaths is a STYLE, not a class: it is loaded over `report'. That base class
%% is not interchangeable — the package defines its abstract, laysummary,
%% acknowledgements and declaration environments in terms of \chapter*, so
%% loading it over `article' fails on the first undefined \chapter.
%%
%% The two University of Edinburgh crest/logo files that the upstream
%% repository carries are NOT redistributable and TeX Live does not install
%% them. This starter therefore bundles no logo and does not reference one —
%% the title page below is the package's own, which is text-only.
\documentclass[11pt,oneside]{report}

%% \title and \author come BEFORE \usepackage{edmaths}, which is the reverse of
%% the usual habit and is not optional here. The package's `hyperref' option
%% bakes \@title and \@author into pdftitle/pdfauthor at load time, so if they
%% are still unset the PDF metadata is filled with LaTeX's own
%% "No \title given" warning text — and the stray braces in that message end
%% the run with "Too many }'s" and no PDF at all.
\title{The Title of Your Thesis}
\author{Your Name}
\date{2026}

%% The degree option is what fills in the title page's degree line, via the
%% package's \degreetext macro. It is REQUIRED: with no degree option the
%% title page literally prints "Specify a degree option!" where the degree
%% should be, which compiles cleanly and ships a broken cover.
%%   phd | masterph | mastersc | mmath | y4project
%%   firstyear | secondyear | thirdyear   (Maxwell Institute annual reports)
%%
%% `onehalf' sets the line spacing (single | onehalf | double). `hyperref' and
%% `fancyhdr' are opt-in here rather than always-on, because the package
%% configures them itself — loading either yourself as well gives you the
%% usual "option clash"/redefinition mess rather than your settings winning.
\usepackage[phd,onehalf,hyperref,fancyhdr]{edmaths}

%% \shorttitle exists only when the `fancyhdr' option is given, and must come
%% AFTER \usepackage{edmaths} — the package defines the macro, so calling it
%% earlier is an undefined control sequence. Without it the running header
%% uses the full title, which wraps badly on a long one.
\shorttitle{Short Title}

%% edmaths loads amsthm but declares NO theorem environments of its own, so
%% these \newtheorem lines are yours to write — without them `theorem' is an
%% undefined environment. Numbering within `chapter' gives Theorem 2.1 rather
%% than a single counter running the length of the thesis.
\theoremstyle{plain}
\newtheorem{theorem}{Theorem}[chapter]
\newtheorem{lemma}[theorem]{Lemma}
\newtheorem{corollary}[theorem]{Corollary}
\theoremstyle{definition}
\newtheorem{definition}[theorem]{Definition}

\begin{document}

\maketitle

%% \declaration takes one argument: pass an EMPTY one to get the School's
%% standard wording (composed by myself, own work, sources acknowledged).
%% Only supply text if your circumstances require a variant — a co-authored
%% chapter, or work previously submitted elsewhere.
\declaration{}

\begin{abstract}
  A single-page summary of the problem, the methods used, the principal
  results and what they mean. Written to be read on its own, so no citations
  and no references to numbered items in the body.
\end{abstract}

%% The lay summary is REQUIRED for Edinburgh PhD theses and is written for a
%% reader outside your field. Delete this environment for a report or an
%% undergraduate project, where it is not asked for.
\begin{laysummary}
  A plain-language account of what the research was about and why it matters,
  aimed at someone with no background in the subject.
\end{laysummary}

\begin{acknowledgements}
  Thank supervisors, collaborators, funders and anyone else who supported the
  work.
\end{acknowledgements}

\dedication{For someone.}

\tableofcontents

\chapter{Introduction}
\label{ch:intro}
Motivate the problem, place it in the existing literature, and set out what
the rest of the thesis does about it.

\chapter{Background}
\label{ch:background}
The definitions and prior results the later chapters rely on.

\begin{theorem}
  \label{thm:example}
  Every bounded sequence in $\mathbb{R}^n$ has a convergent subsequence.
\end{theorem}

\begin{proof}
  Standard; see any text on real analysis.
\end{proof}

%% amsmath, amssymb, amscd, amsthm, geometry, ifthen and setspace are all
%% loaded BY edmaths, so none of them need a \usepackage here. `proof' comes
%% from amsthm and works out of the box; `theorem' does not, which is why the
%% \newtheorem declarations are in the preamble above.

\chapter{Results}
\label{ch:results}
The contribution. Theorem~\ref{thm:example} is used here, and
Table~\ref{tab:example} reports the computations.

\begin{table}[htbp]
  \centering
  \begin{tabular}{lrr}
    \hline
    Case & Iterations & Residual \\
    \hline
    Baseline & 240 & $1.2\times10^{-6}$ \\
    Improved & 118 & $4.7\times10^{-9}$ \\
    \hline
  \end{tabular}
  \caption{Convergence of the two schemes.}
  \label{tab:example}
\end{table}

\chapter{Conclusion}
\label{ch:conclusion}
What was established, what remains open, and where the work goes next.

\appendix

\chapter{Supplementary Material}
Longer computations and data that would interrupt the main argument.

\end{document}
